\documentclass[11pt]{article}

\usepackage{amsmath,amssymb,amsthm,enumitem}

\newcommand{\Mg}{\mathcal M_g}
\newcommand{\Mgn}[2]{\mathcal M_{#1,#2}}
\newcommand{\bMg}{\overline{\mathcal M}_g}
\newcommand{\bMgn}[2]{\overline{\mathcal M}_{#1,#2}}
\newcommand{\DeltaG}{\Delta_{\Gamma}}

\newtheorem{proposition}{Proposition}
\newtheorem{remark}{Remark}

\begin{document}

\title{A counterexample to the proposed boundary-rigidity statement}
\author{}
\date{}
\maketitle

\begin{proposition}
The statement in the question is false under the hypotheses as stated.
\end{proposition}

\begin{proof}
We construct a counterexample in genus \(3\).  The construction uses only standard
facts about cyclic covers and plumbing coordinates near the Deligne--Mumford
boundary.

Consider the two-dimensional Hurwitz locus \(H \subset \mathcal M_3\) consisting
of cyclic triple covers of \(\mathbb P^1\) with five branch points and local
monodromy data
\[
(1,1,1,1,2) \in (\mathbb Z/3\mathbb Z)^5.
\]
Equivalently, after choosing affine coordinates, a dense open part of \(H\) is
given by curves
\[
w^3=\prod_{j=1}^4 (x-a_j)(x-a_5)^2,
\]
with the \(a_j\)'s distinct, modulo \(\mathrm{PGL}_2\).  The Riemann--Hurwitz
formula gives genus \(3\), since each branch point has ramification index \(3\):
\[
2g-2 = 3(-2)+5\cdot 2 =4.
\]
Thus \(g=3\).  The quotient orbifold is a sphere with five order-three orbifold
points, so the corresponding fixed locus in Teichmuller space is naturally
identified with \(\mathcal T_{0,5}\), hence has complex dimension \(2\).

This locus is totally geodesic for the Teichmuller metric.  Indeed, it is the
fixed-point locus of a finite-order mapping class.  If two points in Teichmuller
space are fixed by that mapping class, then the unique Teichmuller geodesic
segment joining them is also fixed, because the mapping class acts by isometries
and preserves the endpoints.  Hence the fixed locus is totally geodesic.  Passing
to the quotient gives a totally geodesic orbifold submanifold of \(\mathcal M_3\).
If desired, one may impose full level structure to remove orbifold issues; the
argument is unchanged.

Now look at the boundary degeneration in which the first two branch points lie
on one component of a nodal rational curve and the remaining three branch points
lie on the other component.  The admissibility condition forces the monodromy at
the node to be \(1\) on the first component and \(2\) on the second component.
Thus the cover over the first component has monodromy type
\[
(1,1,1),
\]
so it is the unique cyclic triple cover of \(\mathbb P^1\) branched at three
points with that monodromy.  This is the equianharmonic elliptic curve
\(Y_0\), with \(j(Y_0)=0\), together with the marked point lying above the node.
The cover over the second component has monodromy type
\[
(1,1,2,2),
\]
and varies in a one-dimensional Hurwitz locus \(C \subset \mathcal M_{2,1}\).
Therefore the corresponding boundary component of \(\overline H\) lies in the
separating boundary stratum
\[
\DeltaG \cong \mathcal M_{1,1}\times \mathcal M_{2,1}
\]
and is
\[
B_0=\{Y_0\}\times C .
\]
In particular, the projection of \(B_0\) to the \(\mathcal M_{1,1}\)-factor is
constant.

Choose a nonconstant holomorphic arc
\[
t \longmapsto Y_t \in \mathcal M_{1,1}
\]
with \(Y_{0}=Y_0\).  For each \(t\), form the family obtained by plumbing the
marked elliptic curve \(Y_t\) to the varying pointed genus-two curve \(Z\in C\).
Concretely, near the stratum
\[
\mathcal M_{1,1}\times \mathcal M_{2,1}\subset \overline{\mathcal M}_3,
\]
the standard plumbing theorem identifies a neighborhood of the boundary with a
neighborhood of the zero section in the normal line bundle whose fiber at
\((Y,Z)\) is
\[
T_{p_Y}Y\otimes T_{p_Z}Z,
\]
where \(p_Y,p_Z\) are the marked points to be glued.  The zero section is the
boundary stratum, and the complement of the zero section parametrizes smooth
genus-three curves.

Let \(D\) be a sufficiently small punctured disk subbundle over \(C\).  For each
\(t\), define \(V_t\subset \mathcal M_3\) to be the image of \(D\) under the
plumbing construction applied to
\[
\{Y_t\}\times C \subset \mathcal M_{1,1}\times \mathcal M_{2,1}.
\]
The construction is holomorphic in \(t\), so the union of the \(V_t\)'s is a
complex submanifold of \(\mathcal M_3\times \Delta\), after possibly shrinking
\(\Delta\).

For \(t=0\), the family just described is precisely a cusp neighborhood in the
Hurwitz surface \(H\).  Hence \(V_0\) is an open complex submanifold of the
totally geodesic surface \(H\).  Under the usual local meaning of ``totally
geodesic,'' this is already enough.  If one insists on a geodesic-convex
interpretation, one may choose \(D\) inside a sufficiently small Teichmuller
horoball for the separating curve being pinched; by the standard log-convexity
of extremal length along Teichmuller geodesics, such a horoball is geodesically
convex inside \(H\).

Now compute the boundary of \(V_t\) in the chosen stratum.  By construction,
\[
\overline{V_t}\cap \DeltaG
=
\{Y_t\}\times C.
\]
Thus
\[
\dim\bigl(\overline{V_t}\cap\DeltaG\bigr)=1
\]
and the projection to the first factor satisfies
\[
\pi_1(\overline{V_t}\cap \DeltaG)=\{Y_t\}.
\]
The family of boundary curves is smoothly trivial, because it is just the fixed
curve \(C\) in the second factor while the point in the first factor moves
holomorphically.

The compatibility condition with smaller strata is also satisfied.  Any further
boundary arises only by allowing the genus-two component \(Z\in C\) to degenerate
inside \(\overline{\mathcal M}_{2,1}\).  For all \(t\), the same degeneration
occurs, with the elliptic factor replaced by \(Y_t\).  Hence the boundary in
smaller strata is obtained from the boundary in the larger stratum in exactly
the same way for every \(t\).

However, by choosing the arc \(t\mapsto Y_t\) nonconstant, we have
\[
\pi_1(\overline{V_t}\cap \DeltaG)=\{Y_t\}\neq \{Y_0\}
\]
for \(t\neq 0\).  Thus the conclusion claimed in the question fails.
\end{proof}

\begin{remark}
The counterexample uses the freedom allowed in the question: the \(V_t\) are
complex submanifolds, not assumed to be closed algebraic subvarieties of
\(\mathcal M_g\).  If one intended an additional global algebraicity or rigidity
hypothesis, it must be stated explicitly; it is not a consequence of the
assumptions given.
\end{remark}

\begin{thebibliography}{9}

\bibitem{HM}
J.~Harris and D.~Mumford,
\emph{On the Kodaira dimension of the moduli space of curves},
Invent. Math. \textbf{67} (1982), 23--88.

\bibitem{ACG}
E.~Arbarello, M.~Cornalba, and P.~Griffiths,
\emph{Geometry of Algebraic Curves. Volume II},
Grundlehren der mathematischen Wissenschaften, vol. 268,
Springer, 2011.

\bibitem{Kerckhoff}
S.~Kerckhoff,
\emph{The asymptotic geometry of Teichmuller space},
Topology \textbf{19} (1980), 23--41.

\end{thebibliography}

\end{document}
```